\title{Direct Preference Optimization: Your Language Model is Secretly a Reward Model}

\begin{abstract}
Large-scale unsupervised language models (LMs) are effective at acquiring extensive world knowledge and some reasoning capabilities; however, directing their outputs with precision remains challenging due to the inherently unsupervised nature of their training process. Traditional approaches for achieving such controllability typically involve collecting human preference data that ranks model outputs, followed by fine-tuning the base LM to better reflect these preferences—commonly using reinforcement learning from human feedback (RLHF). Nevertheless, RLHF tends to be both intricate and potentially unstable, requiring an initial training phase for a reward model that encodes human preferences, after which the LM undergoes reinforcement learning to maximize the predicted reward, all while attempting to preserve fidelity to the original parameters. This work proposes an alternative parameterization of the RLHF reward model that enables direct computation of the associated optimal policy in a closed form, thereby transforming the usual RLHF challenge into a simple classification task. Our proposed algorithm, termed \textit{Direct Preference Optimization} (DPO), offers stability, high performance, and low computational overhead. It removes the need for language model sampling during fine-tuning and obviates extensive hyperparameter search. Experimental results demonstrate that DPO enables language models to be fine-tuned for human alignment as effectively or even more so than prevailing methods. Remarkably, DPO outperforms PPO-based RLHF at sentiment control in generated text, and is competitive or superior for response quality in tasks like summarization and single-turn dialogue, all while being much simpler to deploy and train.
\end{abstract}

\section{Introduction}
Large-scale unsupervised language models (LMs) trained on massive corpora exhibit unexpected abilities~\citep{chowdhery2022palm, brown2020language, touvron2023llama,bubeck2023sparks}. Despite these remarkable competencies, such models are exposed to human-generated data shaped by diverse intentions, backgrounds, and proficiency levels. Not all of these characteristics are desirable for an AI system to replicate; for instance, although it is advantageous for an AI coding assistant to \textit{recognize} typical programming errors so it can suggest corrections, it is preferable for the assistant to emulate the exceptional, if infrequent, high-quality coding examples embedded in its training data during code generation. In a similar vein, a language model should be \textit{informed} about prevalent misconceptions, but we do not intend for it to endorse these misconceptions proportionately to their prevalence, such as repeating a false belief in 50\% of relevant outputs. Effectively filtering and emphasizing the model’s \emph{preferred outputs and conduct} from its broader \textit{range of competencies and knowledge} is therefore fundamental to the development of reliable, high-performing, and governable AI systems \citep{ouyang2022training}. Although contemporary approaches usually employ reinforcement learning (RL) to adjust language models according to human preferences, we demonstrate that the RL-based loss typically applied can in fact be directly optimized via a straightforward binary cross-entropy loss, which streamlines the entire preference alignment process.

At a conceptual level, standard approaches encourage the target behaviors in LMs by leveraging curated datasets encapsulating human preferences for safe and useful responses. This process of preference alignment generally follows the unsupervised pre-training phase, wherein the model learns from vast, unannotated text. While the most direct strategy for learning human preferences consists of supervised fine-tuning on annotated demonstrations of superior responses, the predominant class of solutions is reinforcement learning from human (or AI-generated) feedback (RLHF/RLAIF; \citep{christiano2017deep,bai2022constitutional}). In RLHF, a reward model is first trained using a collection of human preference annotations; subsequently, RL is employed to tune the language model’s policy to favor responses with higher predicted reward while ensuring the policy does not deviate too radically from its pre-trained baseline. Although RLHF yields language models with outstanding dialogue and code generation capabilities, the workflow is significantly more complicated compared to supervised approaches, typically necessitating the training of multiple models and sampling in the RL training loop, which results in considerable computational expenditure.

This work introduces a method for training language models to align with human preferences directly, eliminating the need for explicit reward modeling or reinforcement learning frameworks. We present \textit{{Direct Preference Optimization} (DPO)}, an approach that, while implicitly targeting the same goal as standard RLHF techniques—namely, maximizing expected reward subject to a KL-divergence penalty—remains straightforward to both implement and train. DPO achieves this by enhancing the log odds of preferred responses over less favored ones and, crucially, by utilizing an adaptive, instance-specific importance weighting scheme that averts the collapse observed with a naive probability ratio approach. As with existing algorithms, DPO employs a theoretical model of preference, such as the Bradley-Terry model \cite{bradley1952rankanalysis}, which quantifies the congruence between a reward function and actual human preferences. However, whereas previous strategies employ the preference model to generate a loss used in training a reward model—subsequently optimizing a policy over this learned reward surface—DPO exploits a variable transformation to formulate the preference loss directly in terms of the policy. With access to a corpus of human-labeled preference pairs, DPO can optimize the language model policy directly using a binary cross-entropy loss, yielding the best-fitting policy for an implicit reward derived from preference judgments.

The principal contribution of this paper is Direct Preference Optimization (DPO), a lightweight algorithm for training language models from preference data without resorting to reinforcement learning. Empirical evaluation demonstrates that DPO matches or surpasses current alternatives, including RLHF with PPO, in tasks such as sentiment control, summarization, and conversational modeling, on models containing as many as 6 billion parameters.

\section{Related Work}

Large self-supervised language models demonstrate the ability to solve some problems in a zero-shot setting \citep{radford2019language} or by leveraging few-shot prompts \citep{gpt3,megatron,chowdhery2022palm}. Nevertheless, their efficacy on specific downstream applications and alignment with user intentions are markedly enhanced when the models are further fine-tuned using datasets comprising instructions and human-generated completions \citep{mishra-etal-2022-cross,sanh2022multitask,chung2022scaling,thoppilan2022lamda}. This process, known as `instruction-tuning', allows large language models to more effectively generalize to instructions outside of those observed during tuning, thereby improving general user experience \citep{chung2022scaling}. While instruction tuning has achieved notable results, gathering \textit{relative} human feedback on the quality of responses—rather than absolute expert demonstrations—is often more feasible. Subsequent research has thus adopted the strategy of fine-tuning models using datasets constructed from human preference data, resulting in better performance on tasks such as translation \citep{kreutzer-etal-2018-reliability}, summarization \citep{stiennon2022learning,ziegler2020finetuning}, story-telling \citep{ziegler2020finetuning}, and following instructions \citep{ouyang2022training,ramamurthy2023is}. Typically, these approaches involve first training a neural reward model that predicts user preferences—commonly modeled via frameworks like the Bradley-Terry model \citep{bradley1952rankanalysis}—followed by fine-tuning the language model to maximize expected reward via reinforcement learning, utilizing algorithms such as REINFORCE \citep{williams1992reinforce}, proximal policy optimization (PPO; \cite{schulman2017proximal}), or related methods \citep{ramamurthy2023is}. Another related avenue involves employing instruction-following LLMs, themselves fine-tuned with human feedback, to produce further synthetic datasets reflecting human preferences for qualities like safety or harmlessness \citep{bai2022constitutional}, guided only by broad text-based rubrics as weak supervision. Collectively, these strategies unify developments in both reinforcement learning for language models~\citep{Ranzato2015SequenceLT,paulus2018a,wu2018learning} and general methodologies for preference-based learning from humans \citep{christiano2017deep,kupcsik2018learning}. However, despite the intuitive advantages of relying on human preference data, optimizing large language models via reinforcement learning remains practically complex. The present work offers a theoretically grounded alternative for optimizing relative preference objectives without the use of reinforcement learning.

Beyond language-based scenarios, the task of learning policies from preferences has been investigated within both bandit and reinforcement learning frameworks, resulting in a variety of proposed methods. In particular, contextual bandit learning guided by preferences or rankings of actions, as opposed to explicit reward signals, is referred to as the contextual dueling bandit (CDB; \cite{yue2012karmed,dudik2015contextual}). When absolute rewards are unavailable, theoretical work on CDBs replaces the conventional concept of an optimal policy with the notion of a \textit{von Neumann winner}, which describes a policy that achieves an expected win rate of at least 50\% against \textit{all} alternative policies \citep{dudik2015contextual}. Notably, in the CDB paradigm, preference feedback is received in an online fashion. By contrast, typical human preference learning leverages a static, offline collection of preference-labeled action pairs \citep{yan2022human}. In a related vein, \textit{preference-based RL} (PbRL) involves learning from binary preference comparisons derived from an \textit{unknown} scoring function rather than from rewards \citep{BusaFekete2014,ruiz2023dueling}. A variety of PbRL algorithms have been proposed, some of which can make use of off-policy preference data; however, these methods typically require an explicit intermediate step where the latent scoring (reward) function is estimated and then optimized \citep{jain2013learning,BusaFekete2014,christiano2017deep,sadigh2017active,kupcsik2018learning}. In contrast, our approach introduces a one-stage policy learning method that directly optimizes the policy according to the observed preferences.

\section{Preliminaries}\label{section:prelims}

We briefly summarize the RLHF procedure as outlined in \citeauthor{ziegler2020finetuning} and adopted in subsequent works (\citep{stiennon2022learning, bai2022training, ouyang2022training}). This process generally unfolds in three main steps: (1) supervised fine-tuning (SFT); (2) collection of preference data and training of a reward model; and (3) reinforcement learning-based optimization.

\textbf{SFT}: The RLHF process typically initiates by adapting a pre-trained language model through supervised learning on curated data for a specific downstream task (such as dialogue or summarization), thereby yielding the model $\pi^\text{SFT}$.

\textbf{Reward Modelling Phase}: In the following stage, prompts $x$ are supplied to the SFT model to generate answer pairs $(y_1, y_2) \sim \pi^\text{SFT}(y \mid x)$. These candidate outputs are then assessed by human annotators, who indicate a preference for one response, denoted $y_w\succ y_l \mid x$, where $y_w$ is the selected and $y_l$ is the non-selected response from the pair $(y_1, y_2)$. These human choices are presumed to originate from an underlying but unobserved reward model $r^*(y, x)$. Several modeling frameworks have been applied to capture preferences; a widely used approach is the Bradley-Terry (BT) model \cite{bradley1952rankanalysis} (with generalizations such as the Plackett-Luce ranking models \citep{plackett1975analysis, luce2012individual} being applicable in cases with multiple ranked outputs). According to the BT model, the distribution over human preferences $p^*$ is formalized as follows:

\begin{equation}\label{eq:bradley-terry}
    p^*(y_1\succ y_2 \mid x)=\frac{\exp\left(r^*(x, y_1)\right)}{\exp\left(r^*(x, y_1)\right) + \exp\left(r^*(x, y_2)\right)}.
\end{equation}

Given a fixed collection of preference comparisons $\mathcal{D}=\bigl\{x^{(i)}, y_w^{(i)}, y_l^{(i)}\bigr\}_{i=1}^N$ sampled according to $p^*$, one can introduce a parametric reward estimator $r_{\phi}(x, y)$ and fit its parameters through maximum likelihood estimation. Interpreting this as a binary classification task, the corresponding negative log-likelihood loss function takes the form:

\begin{equation}\label{eq:reward_model}
    \mathcal{L}_R(r_{\phi}, \mathcal{D}) = -\mathbb{E}_{(x, y_w, y_l)\sim \mathcal{D}}\bigl[\log \sigma(r_{\phi}(x, y_w)- r_{\phi}(x, y_l))\bigr]
\end{equation}

where $\sigma$ denotes the logistic sigmoid. For language model settings, the architecture $r_{\phi}(x, y)$ is typically based on the SFT model $\pi^\text{SFT}(y \mid x)$, augmented by appending a linear projection to the last transformer output in order to yield a single scalar score, as described by \cite{ziegler2020finetuning}. To stabilize the reward estimation and control its variability, previous studies normalize the predicted rewards so that $\mathbb{E}_{x,y\sim \mathcal{D}}\left[r_\phi(x, y)\right] = 0$ across all $x$.

\textbf{RL Fine-Tuning Phase}: At the reinforcement learning stage, the acquired reward signal is used to direct updates to the language model. Drawing from established work~\citep{jaques2017sequence, jaques2020human}, the optimization objective is given by:

\begin{equation}\label{eq:RL}
\max_{\pi_{\theta}}  \mathbb{E}_{x\sim \mathcal{D}, y\sim \pi_{\theta}(y \mid x)}\bigl[r_{\phi}(x, y)\bigr] - \beta\mathbb{D}_{\textrm{KL}}\bigl[\pi_{\theta}(y\mid x)\mid \mid \pi_\text{ref}(y\mid x)\bigr],
\end{equation}

with $\beta$ regulating how much the learned policy can diverge from a reference model $\pi_\text{ref}$, typically chosen as the original SFT policy $\pi^\text{SFT}$. The language model $\pi_\theta$ is also seeded with the SFT initialization. This KL regularization is essential, as it constrains the policy close to regions where the reward model is reliable, preserves the diversity of generations, and avoids pathological concentration on a few high-reward outputs. Since generation is over discrete tokens, this objective cannot be directly differentiated and thus calls for reinforcement learning methods for optimization. The prevailing technique \citep{ziegler2020finetuning, stiennon2022learning, bai2022training, ouyang2022training} reformulates the reward as ${r(x, y) = r_{\phi}(x, y) -\beta (\log \pi_{\theta}(y\mid x) - \log \pi_\text{ref}(y\mid x))}$ and applies PPO \cite{schulman2017proximal} to maximize it.

\section{Direct Preference Optimization}\label{sec:DPO}

Recognizing the complexities of employing RL techniques for large-scale language model adaptation, our objective is to establish a streamlined method for policy optimization directly informed by preferences. In contrast to conventional RLHF pipelines, which separate reward learning and RL-based policy optimization, our framework utilizes a specific reward model structure that permits analytic determination of the associated optimal policy, thus circumventing the need for reinforcement learning loops. As elaborated in the following, the main insight is to utilize a closed-form correspondence between reward models and their optimal induced policies. This relationship allows us to convert a reward-focused loss into a policy-oriented loss through a variable substitution, obviating the need for explicit standalone reward network fitting while remaining consistent with widely used preference models, such as the Bradley-Terry framework. Consequently, the policy model implicitly serves as both the language generator and as a representation of the reward function.

\textbf{Derivation of the DPO Objective.} We begin from the reinforcement learning objective defined in Eq.~\ref{eq:RL} with a general reward function $r$, as considered in previous studies. Prior works~\citep{peters2007reinforcement, peng2019advantage, korbak2022reinforcement, go2023aligning} have established that the optimal policy for the KL-regularized reward maximization objective in Eq.~\ref{eq:RL} is:

\begin{equation}\label{eq:op_policy}
    \pi_r(y\mid x) = \frac{1}{Z(x)}\pi_\text{ref}(y\mid x)\exp\left(\frac{1}{\beta}r(x, y)\right),
\end{equation}

where $Z(x) =\sum_{y}\pi_\text{ref}(y\mid x)\exp\left(\frac{1}{\beta}r(x, y)\right)$ represents the partition function. For detailed steps, refer to Appendix \ref{app:derivation1}. Even assuming access to the MLE-based reward estimate $r_{\phi}$ as a proxy for the true reward $r^*$, computing $Z(x)$ remains computationally demanding~\citep{korbak2022reinforcement, go2023aligning}, which hampers practical application of this form. However, Eq.~\ref{eq:op_policy} can be algebraically manipulated to isolate the reward function in terms of its optimal policy $\pi_r$, the reference distribution $\pi_\text{ref}$, and the partition function $Z(\cdot)$. Specifically, by applying the logarithm to both sides of Eq.~\ref{eq:op_policy} and simplifying, we derive:

\begin{equation}\label{eq:main_eq}
    r(x,y) =\beta \log \frac{\pi_r(y\mid x)}{\pi_\text{ref}(y\mid x)} + \beta \log Z(x).
\end{equation}

This change of variables can be applied for the ground-truth reward $r^*$ and the associated optimal policy $\pi^*$. Importantly, the Bradley-Terry framework utilizes only the difference between reward values for a pair of outputs, i.e., ${p^*(y_1 \succ y_2 \mid x) = \sigma(r^*(x, y_1) - r^*(x, y_2))}$. Inserting the parameterization from Eq.~\ref{eq:main_eq} for $r^*(x, y)$ into the Bradley-Terry preference formula in Eq.~\ref{eq:bradley-terry} causes the partition function to cancel, yielding a preference probability depending solely on the optimal policy $\pi^*$ and the reference policy $\pi_\text{ref}$. Consequently, the optimal RLHF policy $\pi^*$ must satisfy:

\begin{equation}\label{eq:objective}
    p^*(y_1\succ y_2 \mid x)=\frac{1}{1 + \exp\left(\beta \log \frac{\pi^*(y_2\mid x)}{\pi_\text{ref}(y_2\mid x)} - \beta \log \frac{\pi^*(y_1\mid x)}{\pi_\text{ref}(y_1\mid x)}\right)}
\end{equation}

A full proof is given in Appendix~\ref{app:derivation2}. Although Eq.~\ref{eq:objective} utilizes the Bradley-Terry model, analogous formulations extend to the broader Plackett-Luce family~\citep{plackett1975analysis, luce2012individual}, as discussed in Appendix~\ref{app:plackett_luce_models}.

Having re-expressed human preference likelihoods in terms of an optimal policy rather than an explicit reward, we can now define a maximum likelihood objective with respect to a parameterized policy $\pi_\theta$. In parallel to the standard reward modeling setup (see Eq.~\ref{eq:reward_model}), this policy objective is given by:

\begin{equation}\label{eq:optimum_model}
    \mathcal{L}_\text{DPO}(\pi_{\theta}; \pi_\text{ref}) = -\mathbb{E}_{(x, y_w, y_l)\sim \mathcal{D}}\left[\log \sigma \left(\beta \log \frac{\pi_{\theta}(y_w\mid x)}{\pi_\text{ref}(y_w\mid x)} - \beta \log \frac{\pi_{\theta}(y_l\mid x)}{\pi

By employing this approach, we implicitly learn a reward function through a different parameterization, wherein the resulting optimal policy directly corresponds to $\pi_\theta$. Furthermore, as our methodology aligns with optimizing a reparametrized Bradley-Terry model, it inherits well-established theoretical guarantees, notably consistency under suitable conditions on the distribution of preference data~\cite{bong2022generalized}. We provide a deeper theoretical analysis of DPO’s properties and its connections to prior research in Section~\ref{sec:theory}.

\textbf{Mechanistic Interpretation of the DPO Update.} To gain mechanistic insight into DPO, we examine the gradient of the loss function $\mathcal{L}_\text{DPO}$. The parameter gradient can be expressed as:

\begin{multline*}\label{eq:gradient}
    \nabla_\theta \mathcal{L}_\text{DPO}(\pi_\theta;\pi_\text{ref}) = \\ -\beta\mathbb{E}_{(x, y_w, y_l) \sim \mathcal{D}} \bigg[\underbrace{\sigma(\hat{r}_\theta(x, y_l) - \hat{r}_\theta (x, y_w))}_\text{higher weight when reward estimate is wrong}\bigg[\underbrace{\nabla_\theta\log \pi(y_w \mid x)}_\text{increase likelihood of $y_w$} - \underbrace{\nabla_\theta\log\pi(y_l \mid x)}_\text{decrease likelihood of $y_l$}\bigg]\bigg],
\end{multline*}

where $\hat{r}_\theta(x, y) = \beta \log \frac{\pi_\theta(y \mid x)}{\pi_\text{ref}(y \mid x)}$ represents the reward function implicitly defined by the current language model $\pi_\theta$ and the reference model $\pi_\text{ref}$ (further elaborated in Section~\ref{sec:theory}). Conceptually, taking the gradient of $\mathcal{L}_\text{DPO}$ results in raising the probability of favored completions $y_w$ while reducing that of unfavored completions $y_l$. Crucially, each example is scaled according to how much the implicit reward model $\hat{r}_\theta$ overvalues the dispreferred response relative to the preferred one, modulated by the parameter $\beta$. This mechanism reflects the degree of misordering by the reward estimator and integrates the influence of the KL constraint's strength. Our empirical results indicate that this weighting plays a critical role; omitting the weighting coefficient, as in a simplistic variant, may lead to undesirable degeneration of the language model (see Appendix Table~\ref{tab:unlikelihood_generations}).

\textbf{DPO framework.} The standard workflow for DPO consists of the following steps: First, completions $y_1, y_2 \sim \pi_\text{ref}(\cdot \mid x)$ are sampled for each prompt $x$, and then human annotators provide preference labels to create an offline preference dataset $\mathcal{D} = \{x^{(i)}, y_w^{(i)}, y_l)^{(i)}\}_{i=1}^N$. Second, the language model $\pi_\theta$ is trained by minimizing $\mathcal{L}_\text{DPO}$ with respect to the specified $\pi_\text{ref}$, dataset $\mathcal{D}$, and the target $\beta$. 

In practice, researchers prefer to utilize pre-existing public preference datasets rather than produce new completions and manually annotate preferences. Since these preference datasets are generally generated using $\pi^\text{SFT}$, it is standard to initialize $\pi_\text{ref}$ as $\pi^\text{SFT}$ if it is accessible. If $\pi^\text{SFT}$ is unavailable, $\pi_\text{ref}$ is initialized by maximizing the likelihood over the preferred responses ${(x, y_w)}$: ${\pi_\text{ref} = \argmax_{\pi}\mathbb{E}_{x, y_w \sim \mathcal{D}}\left[\log \pi(y_w \mid x)\right]}$. This strategy aims to reduce the distributional mismatch between the unknown true reference distribution and the $\pi_\text{ref}$ employed during DPO training. Comprehensive information on implementation specifics and selected hyperparameters can be found in Appendix~\ref{app:implementation}.

\section{Theoretical Analysis of DPO} In what follows, we present a deeper conceptual understanding of DPO, establish its theoretical foundation, and discuss its benefits in light of shortcomings commonly encountered in actor-critic approaches for RLHF, such as PPO~\cite{schulman2017proximal}.

\label{sec:theory}

\subsection{Your Language Model Is Secretly a Reward Model} DPO forgoes both explicit reward modeling and reinforcement learning for policy improvement, opting instead for a unified maximum likelihood framework. In particular, the objective in Eq. \ref{eq:main_eq} can be interpreted as a Bradley-Terry model where rewards are given by $r^*(x, y) = \beta \log\frac{\pi^*_\theta(y \mid x)}{\pi_\text{ref}(y \mid x)}$. Here, we train the model $\pi_\theta$ analogously to how reward models are trained via Eq. \ref{eq:reward_model}, under an appropriate variable substitution. In this section, we construct the formal theory for this parameterization, show that it imposes no additional limitations on the expressivity of learned reward models, and that it enables direct recovery of the optimal policy. We start by establishing an equivalence relation among reward functions.

\begin{definition}
We define two reward functions $r(x, y)$ and $r'(x, y)$ as equivalent if and only if there exists a function $f$ such that ${r(x, y)-r'(x, y) = f(x)}$.
\end{definition}

This equivalence is readily seen to establish an equivalence relation, thus dividing the space of reward functions into distinct equivalence classes. Based on this, we have the following lemmas:

\begin{lemma}\label{lemma:same_prefrence} Within the Plackett-Luce, and specifically the Bradley-Terry, preference frameworks, any pair of reward functions belonging to the same equivalence class result in identical preference distributions.
\end{lemma}

\begin{lemma}\label{lemma:same_policy}
    Any two reward functions within the same equivalence class induce identical optimal policies under the specified constrained RL setting.
\end{lemma}

The arguments for these results are direct and are presented in Appendix \ref{app:lemma1}. The first lemma reflects a widely recognized identifiability issue inherent in the Plackett-Luce model family \cite{plackett1975analysis}. Due to this ambiguity, it is standard practice to introduce additional constraints to ensure meaningful MLE recovery in Eq.~\ref{eq:reward_model} \cite{bong2022generalized}. The second lemma confirms that every member of a reward equivalence class specifies the same optimal policy; therefore, our learning goal is restricted to recovering a representative reward from the relevant class. The following theorem, formally proven in Appendix~\ref{app:thm1}, characterizes these equivalence classes:

\begin{theorem}\label{thm:main}
    Assuming mild regularity conditions, every reward equivalence class compatible with Plackett-Luce (and in particular, Bradley-Terry) models can be expressed in the form ${r(x, y) = \beta \log \frac{\pi(y\mid x)}{\pi_\text{ref}(y\mid x)}}$ for some probabilistic model $\pi(y\mid x)$ and a given reference model $\pi_\text{ref}(y \mid x)$.
\end{theorem}

\begin{sproof}
    Take an arbitrary reward function $r(x, y)$, and let it induce the optimal policy $\pi_r(y \mid x)$, as defined in Eq. \ref{eq:op_policy}. We will show that there exists a representative of the equivalence class of $r$ with the proposed reparameterized structure. Consider the transformation 
\begin{equation}
    f(r; \pi_\text{ref}, \beta)(x, y) = r(x, y) - \beta\log\sum_{y}\pi_\text{ref}(y\mid x)\exp\left(\frac{1}{\beta}r(x, y)\right)
\end{equation}
Here, $f$ adjusts the reward function by subtracting the log-partition term corresponding to $\pi_r$, which only depends on $x$. Hence, $f(r; \pi_\text{ref}, \beta)(x, y)$ remains within the equivalence class of $r(x, y)$. Now, if we set $r$ to the expression on the right-hand side of Eq.~\ref{eq:main_eq} (valid for any $r$), the mapping yields $f(r; \pi_\text{ref}, \beta)(x, y) = \beta \log \frac{\pi_r(y\mid x)}{\pi_\text{ref}(y\mid x)}$. Thus, this projection $f$ returns a form-equivalent reward function within the same equivalence class, so the proposed parameterization suffices to cover all cases for the reward model.
\end{sproof}

Alternatively, Theorem~\ref{thm:main} can be interpreted as identifying the specific reward function, among all those equivalent up to a constant shift, that the DPO reparameterization chooses. In particular, it selects the reward function that fulfills the following condition:

\begin{equation}\label{eq:lag_p}
     \sum_{y}\underbrace{\pi_\text{ref}(y\mid x)\exp\left(\frac{1}{\beta}r(x, y)\right)}_{=\pi(y\mid x)\text{, using Thm.~\ref{thm:main} reparam.}} = 1,
\end{equation}

This guarantees that $\pi(y\mid x)$ defines a legitimate probability distribution, with all probabilities non-negative and normalized to sum to 1. Referring back to Eq.~\ref{eq:op_policy}, we observe that Eq.~\ref{eq:lag_p} represents the partition function corresponding to the optimal policy derived from the reward $r(x, y)$. The fundamental insight behind the DPO algorithm lies in strategically constraining the otherwise under-specified Plackett-Luce class (including, notably, the Bradley-Terry model) of preference models. By doing so, one preserves the expressive capacity of possible reward models while ensuring that the optimal policy from Eq.~\ref{eq:op_policy} can always be computed in closed form for any prompt $x$.

\subsection{Instability of Actor-Critic Algorithms} Our framework further sheds light on sources of instability in standard actor-critic approaches used in RLHF, such as PPO. Adhering to the RLHF pipeline, we restrict our attention to the reinforcement learning fine-tuning step described in Section~\ref{section:prelims}. Connections emerge with the control-as-inference framework \cite{levine2018reinforcement} for constrained RL formulated in \ref{eq:RL}. Given a parametrized policy $\pi_{\theta}(y\mid x)$, optimization proceeds by minimizing $\mathbb{D}_{\text{KL}}[\pi_{\theta}(y|x) \mid \mid \pi^*(y\mid x)]$, where $\pi^*$ denotes the optimal policy determined in Eq.~\ref{eq:optimum_model} based on the reward function $r_{\phi}(y, x)$. After manipulation, this yields the following objective:

\begin{equation}\label{eq:AC}
    \max_{\pi_{\theta}}\mathbb{E}_{\pi_{\theta}(y\mid x)}\bigg[\underbrace{r_{\phi}(x, y) -\beta\log\sum_{y}\pi_\text{ref}(y\mid x)\exp\left(\frac{1}{\beta}r_{\phi}(x, y)\right)}_{f(r_{\phi}, \pi_\text{ref}, \beta)} - \underbrace{\beta\log\frac{\pi_{\theta}(y\mid x)}{\pi_\text{ref}(y\mid x)}}_{\text{KL}}\bigg]
\end{equation}

The objective here is identical to that used in previous studies 
\citep{ziegler2020finetuning, stiennon2022learning, bai2022training, ouyang2022training}, employing the DPO-equivalent reward for the reward class $r_{\phi}$. Within this context, the normalization factor in $f(r_{\phi}, \pi_\text{ref}, \beta)$ can be understood as the soft value function corresponding to the reference policy $\pi_\text{ref}$. Although this normalization does not influence the optimal policy itself, omitting it may cause the policy gradient estimator to exhibit significant variance, potentially leading to unstable optimization dynamics. This normalization can be estimated via a learned value function, but that approach may also introduce optimization challenges. As an alternative, previous work often employed normalization by referencing a baseline from human completions, which serves as a single-sample Monte Carlo approximation of the normalization constant. In contrast, the DPO formulation inherently constructs a reward function that obviates the need for such baselines.

\section{Experiments}
Here, we present empirical analyses of DPO's capability to directly optimize policies from preference data. To begin, in a controlled environment for text generation, we investigate how well DPO balances maximizing the reward with minimizing KL-divergence from the reference policy, especially in comparison to typical preference optimization algorithms like PPO. We then explore DPO's effectiveness on larger-scale models and more challenging RLHF problems, such as summarization and conversational tasks. Notably, we observe that DPO, with minimal hyperparameter adjustment, matches or surpasses established approaches like RLHF using PPO, and even outperforms simply selecting the top trajectory out of $N$ samples based on a learned reward. We detail our experimental methodology below, with supplementary information provided in Appendix~\ref{app:exp_details}.

\textbf{Tasks.} Our study investigates three distinct tasks involving open-ended text generation. Across all tasks, algorithms are trained to optimize a policy using a dataset of preference samples, denoted as $\mathcal{D}=\bigl\{x^{(i)}, y_w^{(i)}, y_l^{(i)}\bigr\}_{i=1}^N$. In the task of \textbf{controlled sentiment generation}, $x$ serves as the opening segment of a movie review sourced from the IMDb dataset \cite{maas-EtAl:2011:ACL-HLT2011}, and the model’s objective is to produce a continuation $y$ that expresses positive sentiment. For systematic evaluation, preference pairs between generated samples are constructed via a pre-trained sentiment classifier, ensuring that $p(\text{positive}\mid x,y_w)>p(\text{positive}\mid x,y_l)$. For supervised fine-tuning (SFT), GPT-2-large is further trained to convergence using the IMDb dataset’s training split (see App~\ref{app:sentiment_details} for implementation details). In the \textbf{summarization} task, $x$ represents a Reddit forum post, and the aim is to generate a concise summary $y$ that captures the post’s central points. Adopting established methodologies, we use the Reddit TL;DR summarization corpus \citep{volske-etal-2017-tl} along with human preference annotations collected by \citeauthor{stiennon2022learning}. Our SFT baseline is fine-tuned on summaries written by humans\footnote{\url{https://huggingface.co/CarperAI/openai_summarize_tldr_sft}} using the TRLX \citep{leandro_von_werra_2023_7790115} RLHF library. The dataset of human preferences was assembled by \citeauthor{stiennon2022learning} utilizing generations from a different, yet similarly-trained, SFT system. Lastly, in the case of \textbf{single-turn dialogue}, $x$ is a user-initiated query, ranging from scientific questions to requests for advice. The model must craft a response $y$ that is both informative and engaging; for this, we leverage the Anthropic Helpful and Harmless dialogue dataset \citep{bai2022training}, which consists of 170,000 dialogues between a person and an AI assistant. Each conversation concludes with two model-generated replies, accompanied by a label indicating which response a human preferred. In this task, as there is no publicly-available pre-trained SFT checkpoint, we fine-tune an existing language model exclusively on preferred response completions to establish the SFT baseline.

\textbf{Evaluation.} We adopt two primary strategies for evaluating our models. For assessing how well each algorithm optimizes the constrained reward maximization objective in the controlled sentiment generation scenario, we examine the trade-off frontier between the achieved reward and KL-divergence from the reference policy. This analysis is feasible since the actual reward function—a sentiment classifier—is accessible. Conversely, in practical applications where the underlying reward function is unknown, we measure algorithmic performance using \textit{win rate} against a baseline policy. In these cases, GPT-4 serves as a stand-in for human evaluators to judge the quality of summaries and the helpfulness of responses in summarization and single-turn dialogue tasks, respectively. For the summarization task, baseline comparisons use reference summaries provided in the test set, while for dialogue, the baseline is defined by the preferred response present in the test data. While prior research has suggested that language models can surpass existing automated metrics in evaluation quality \citep{Chen2023ExploringTU}, we further validate the suitability of GPT-4 for evaluation by running a human study detailed in Sec.~\ref{sec:human-judgments}. Results indicate a strong alignment between GPT-4 and human judgments, with human raters' agreement with GPT-4 being at least comparable to, and often exceeding, agreement rates between different human annotators.

\textbf{Methods.} Beyond DPO, we benchmark several other approaches for training language models to capture human preferences. As baselines, we utilize zero-shot prompting with \textbf{GPT-J} \citep{gpt-j} for summarization and 2-shot prompting with \textbf{Pythia-2.8B} \citep{biderman2023pythia} in dialogue. We include the \textbf{SFT} model and the \textbf{Preferred-FT} model, which is obtained by fine-tuning with supervised learning on preferred completions $y_w$, either produced by the SFT model (for sentiment control and summarization) or by a standard language model (for dialogue). We also evaluate a pseudo-supervised technique, \textbf{Unlikelihood}~\citep{welleck2019neural}, which trains the policy to maximize probability for $y_w$ while explicitly minimizing probability for $y_l$; here, an optional weighting coefficient $\alpha\in[0,1]$ is applied to the unlikelihood loss term. Additional methods under consideration include \textbf{PPO} \citep{schulman2017proximal} employing a reward model trained from preference data, as well as \textbf{PPO-GT}, which utilizes the oracle ground-truth reward in the controlled sentiment case. For sentiment experiments, PPO-GT is implemented both with a standard version \cite{leandro_von_werra_2023_7790115} and a variant with reward normalization and further hyperparameter tuning to enhance performance; analogous adjustments are also applied to PPO with learned rewards. Finally, we incorporate the \textbf{Best of $N$} approach, in which $N$ outputs are sampled from the SFT model (or from Preferred-FT for dialogue), and the sample with the highest reward—according to a learned reward model—is selected. While this approach yields strong performance by isolating reward model quality from PPO optimization challenges, it is computationally expensive for larger $N$, as it entails sampling $N$ responses for every test query.

\subsection{How well can DPO optimize the RLHF objective?}

The standard KL-constrained reward maximization objective implemented in most RLHF methods seeks to maximize the reward while simultaneously penalizing excessive divergence from a reference policy. As such, fair algorithmic comparisons necessitate evaluating both the attained reward and the extent of KL divergence—gains in reward at the cost of substantially higher KL are not necessarily advantageous. In Figure~\ref{fig:frontier-tldr-main}, we illustrate the trade-off between reward and KL for a variety of methods within the sentiment classification context. For each algorithm, we conduct multiple experiments, each with a different policy conservativeness hyperparameter (target KL $\in\{3,6,9,12\}$ for PPO, $\beta \in \{0.05,0.1,1,5\}$, $\alpha\in\{0.05,0.1,0.5,1\}$ for unlikelihood, and varying random seeds for preferred-FT). This process covers a total of 22 experimental runs. Following every 100 training updates up to convergence, we assess each learned policy across a suite of test prompts, computing the mean reward according to the actual reward function as well as the mean sequence-level KL\footnote{Calculated as the sum of KL-divergences at each generation step.} with the reference policy, $\text{KL}\left(\pi\mid \mid \pi_\text{ref}\right)$. Our empirical findings indicate that DPO establishes a substantially superior frontier, simultaneously attaining higher reward and maintaining low KL. This outcome is particularly striking for two key reasons. First, despite both DPO and PPO targeting the identical objective, DPO demonstrates clear efficiency gains; the reward/KL frontier of DPO consistently outperforms that of PPO. Second, DPO surpasses PPO's frontier, even in scenarios where PPO leverages oracle reward information (PPO-GT).

\subsection{Can DPO scale to real preference datasets?}
\label{sec:dpo-real-datasets}
We proceed to assess the effectiveness of DPO fine-tuning on tasks involving summarization and single-turn dialogue. For the task of summarization, traditional automated metrics like ROUGE often show weak alignment with actual human preferences~\citep{stiennon2022learning}. Previous studies have indicated that language models fine-tuned with PPO on human feedback tend to generate higher-quality summaries. Our evaluation involves sampling outputs on the TL;DR summarization dataset's test split and measuring the mean win rate when compared against the corresponding reference summaries from the test set. To ensure fair comparison across methods, all models generate completions with sampling temperatures ranging from 0.0 to 1.0; win rates for each temperature are presented in Figure~\ref{fig:frontier-tldr-main} (right). DPO, PPO, and Preferred-FT each fine-tune the identical GPT-J SFT checkpoint\footnote{\url{https://huggingface.co/CarperAI/openai_summarize_tldr_sft}}. Our findings indicate that at temperature 0.0, DPO achieves an average win rate of roughly 61\%, surpassing PPO's best performance of about 57\% under the same temperature conditions. Furthermore, DPO attains a higher maximum win rate than the best-of-$N$ baseline approach. It is important to highlight that the $\beta$ hyperparameter for DPO was not extensively optimized, which suggests these results could be further improved. Notably, DPO demonstrates superior stability across various temperatures, whereas PPO’s performance can drop to that of the original GPT-J model when the sampling temperature increases. In contrast, Preferred-FT yields only marginal gains over the initial SFT model. A direct comparison between DPO and PPO using human evaluators is detailed in Section~\ref{sec:human-judgments}, with DPO outputs sampled at temperature 0.25 being favored 58\% of the time over PPO outputs at the same temperature.

For single-turn dialogue, our assessment involves applying each method to the portion of the Anthropic HH dataset \citep{bai2022training} test split that features a single human-assistant exchange. To evaluate with GPT-4, we use the preferred completions in the test set as the gold standard and calculate win rates for each approach accordingly. Since a standard SFT model is not established for this task, we initialize with the pre-trained Pythia-2.8B, then employ Preferred-FT to train a reference model on selected completions to ensure outputs remain in-distribution, and subsequently proceed to DPO training. Our experiments also include comparisons with the best sample out of 128 completions produced by Preferred-FT (noting that the Best of $N$ baseline reaches a performance ceiling at $N=128$ for this evaluation; further details appear in Appendix Figure~\ref{fig:best-of-n}) and a 2-shot prompted variant of the Pythia-2.8B base model. The findings reveal that DPO matches or surpasses other methods at their optimal temperature settings. Additionally, we analyze an RLHF model trained using PPO on the Anthropic HH dataset \footnote{\url{https://huggingface.co/reciprocate/ppo_hh_pythia-6B}}, referenced from a reputable source \footnote{\url{https://github.com/CarperAI/trlx/tree/main/examples/hh}}, but are unable to identify a prompt or sampling temperature that allows this model to outperform the base Pythia-2.8B. Considering our TL;DR results and the fact that both PPO and Best of 128 target the same reward objective, we treat the Best of 128 metric as an approximate benchmark for PPO performance. In summary, DPO stands out as the only method that is both computationally efficient and achieves improvements beyond the Anthropic HH dataset's preferred completions, delivering performance on par with or superior to the resource-intensive Best of 128 strategy. Figure~\ref{fig:dialogue-main} further demonstrates that DPO attains peak performance within a comparatively small number of training steps.

\subsection{Generalization to a new input distribution}

To investigate how well PPO and DPO policies adapt to unseen data, we assess their performance from the Reddit TL;DR summarization task on a different distribution: the test set of CNN/DailyMail news articles \citep{nallapati-etal-2016-abstractive}. We employ the top-performing sampling temperatures from the TL;DR evaluations (0 and 0.25) for this assessment. Table~\ref{tab:ood} reports the outcomes. For evaluation, we determined the GPT-4 win rate by comparing model-generated summaries to the gold-standard CNN/DailyMail references, using the same GPT-4 (C) prompt as in Reddit TL;DR experiments, except substituting ``forum post'' with ``news article''. Notably, on this novel distribution, DPO retains a clear advantage over the PPO policy. These findings suggest that DPO exhibits comparable, if not superior, generalization capabilities relative to PPO, despite DPO lacking access to the additional unlabeled Reddit TL;DR prompts utilized by PPO.

\subsection{Validating GPT-4 judgments with human judgments}
\label{sec:human-judgments}
To assess the trustworthiness of GPT-4’s evaluations, we performed a human subject study using outputs from the TL;DR summarization experiment and tested two GPT-4 prompting strategies. The \textbf{GPT-4 (S)} (simple) prompt asks which summary better captures the key information from the post. In contrast, the \textbf{GPT-4 (C)} (concise) prompt additionally queries which summary is more concise, as we noticed that under the \textbf{GPT-4 (S)} prompt, GPT-4 tends to favor longer, more repetitive summaries compared to human preferences. Complete prompt texts are available in Appendix~\ref{app:prompts}. We conducted three comparative evaluations: the best-performing (DPO, temp. 0.25), the worst-performing (PPO, temp. 1.0), and an intermediate-performing (SFT, temp. 0.25) model, each pitted against the PPO policy with its best-performing (greedy) sampling temperature. This strategy samples a broad range of summary quality; multiple human raters were recruited for the DPO and PPO-1 matchups. Our results show that GPT-4’s assessments align with human preferences as frequently as humans agree amongst themselves, for both prompting styles. This supports using GPT-4 as a viable substitute for human evaluation, at least within our experimental constraints. Moreover, the \textbf{GPT-4 (C)} prompt yields win rates more closely tracking those of human evaluators; thus, we adopt this prompt as our primary evaluation method in Section~\ref{sec:dpo-real-datasets}. Further information about the human subject study, including screenshots of the web interface and the roster of raters, can be found in Appendix~\ref{app:human-study}.

\section{Discussion}
Preference-based learning presents a robust and scalable approach for developing language models that are both capable and aligned with human values. In this work, we have presented DPO—a straightforward method for training language models from preferences, which operates independently of reinforcement learning. Instead of transforming the preference learning challenge into a conventional RL problem to leverage established RL techniques, DPO establishes a correspondence between language model policies and reward functions, permitting direct optimization of language models to reflect human preferences using a standard cross-entropy loss. This eliminates the necessity for reinforcement learning and does not sacrifice generality. With minimal hyperparameter adjustment, DPO achieves comparable or superior results relative to prevailing RLHF techniques, such as those using PPO. Consequently, DPO substantially lowers the barriers associated with preference-based training of language models.

\textbf{Limitations \& Future Work.} Our findings highlight several avenues for future investigation. For instance, how does generalization out of distribution for DPO-trained policies compare to approaches utilizing explicit reward functions? Preliminary evidence indicates that DPO models exhibit generalization performance on par with those trained using PPO, but a more exhaustive analysis is required. It also remains to be seen whether self-labeling through the DPO policy can facilitate effective utilization of prompts lacking labels. Another open question concerns the dynamics of reward over-optimization within the direct preference optimization framework, particularly in relation to the modest drop in performance shown in Figure~\ref{fig:dialogue-main}-right. Furthermore, while our experiments have involved models up to 6B parameters, extending DPO to significantly larger, state-of-the-art models represents a promising direction for subsequent research. On the topic of evaluation, we observe that GPT-4's win rates are influenced by the specific prompt formulation; further work could investigate optimal strategies for obtaining reliable assessments from automated evaluators. Finally, DPO's potential applications extend beyond language model alignment, offering prospects for training generative models across diverse domains.